\section{Where the idea came from: codes in the brain}
\label{sec:neural}

This section is the origin story. It is not a prerequisite for the technical
content --- a reader in a hurry can skip to Section~\ref{sec:fca} --- but it
explains \emph{why} anyone would build a system like \odag{} in the first
place, and it contains the observation that the rest of the paper formalises.

\subsection{What a neural code is}

A sensory neuron does not report a number to a central reader. It fires, or it
does not, and what it means is determined entirely by which things in the
world make it fire. Record from many neurons at once while showing many
stimuli, threshold each response into ``active'' or ``silent'', and you have a
matrix:

\begin{center}
\begin{tikzpicture}[font=\small]
\matrix (m) [matrix of nodes, nodes={draw, minimum size=6mm, anchor=center,
             inner sep=0pt, font=\small}, column sep=-\pgflinewidth,
             row sep=-\pgflinewidth,
             row 1/.style={nodes={draw=none, fill=none}},
             column 1/.style={nodes={draw=none, fill=none, text width=2.1cm,
                              align=right}}]
{
              & $n_1$ & $n_2$ & $n_3$ & $n_4$ & $n_5$ & $n_6$ \\
  face A      & 1     & 1     & 0     & 1     & 0     & 0     \\
  face B      & 1     & 1     & 0     & 0     & 1     & 0     \\
  cup         & 0     & 0     & 1     & 0     & 0     & 1     \\
  hand        & 0     & 1     & 1     & 0     & 0     & 1     \\
};
\node[font=\footnotesize\itshape, text=odgrey, anchor=west, align=left]
  at ([xshift=6mm]m.east)
  {rows: stimuli\\columns: neurons\\cell: did it fire?};
\end{tikzpicture}
\end{center}

\noindent
Real cortical codes of this kind are \emph{sparse}: for any one stimulus only
a small fraction of the units are active \citep{barlow1972,foldiak1990,%
olshausen1996}. Sparseness is not an accident of measurement. It is what you
get if you ask a network to represent its input economically --- to spend few
active units per item while keeping items distinguishable --- and it is the
representational regime in which the observation below has teeth.

Between two extremes:

\begin{itemize}[leftmargin=1.4em]
\item A \textbf{localist} code gives each thing its own unit --- the
  ``grandmother cell''. Simple to read, but you need as many units as things,
  and nothing is shared between related items.
\item A \textbf{dense distributed} code has roughly half the units active for
  everything. Maximally efficient in bits, but the representation of one item
  overlaps every other for no meaningful reason.
\item A \textbf{sparse} code sits between: each item activates a small set,
  and --- crucially --- \emph{related items activate overlapping sets}.
\end{itemize}

That last clause is where structure appears.

\subsection{The observation}

Take the code for ``spaniel'' and the code for ``dog''. In a sparse code that
has learned anything about the world, the units active for a spaniel will
include the units that respond to dogs in general, plus some extra ones
specific to spaniels. Write $\code{x}$ for the set of active units of
stimulus $x$. Then
\[
  \code{\text{dog}} \;\subseteq\; \code{\text{spaniel}} .
\]

Read that line slowly, because it is the whole idea. The taxonomic statement
``a spaniel is a dog'' has become a containment between two sets of active
units. \textbf{More specific means more bits on.} And containment brings its
entire order theory (Section~\ref{sec:orders}) along with it: the codes are
partially ordered by $\subseteq$, incomparable things are exactly things
neither of which is a special case of the other, and the covering relation of
that order is a Hasse diagram --- a category graph.

\begin{keyidea}
\textbf{The category graph need not be stored anywhere.} No synapse has to
encode the arrow from spaniel to dog. The arrow is a \emph{consequence} of
which units overlap. Associations between concepts can be implicit in the
code rather than carried by connections \citep{foldiak2003}.
\end{keyidea}

This has three immediate consequences that shaped \odag{}'s design, and each
of them will reappear later in the paper wearing engineering clothes.

\begin{enumerate}[leftmargin=1.6em]
\item \textbf{Queries are set operations.} ``Find everything that is both a
  dog and brown'' becomes: find every $x$ whose code contains the dog units
  \emph{and} the brown units. A conjunctive query is an intersection --- or,
  in hardware terms, a bitwise \textsc{and}. This is \odag{}'s only query
  primitive (\S\ref{sec:get}).
\item \textbf{Similarity is overlap.} Two items with largely overlapping codes
  are similar without anyone having said so. Graded retrieval --- ``what is
  most like this?'' --- is $|\code{x} \cap \code{y}|$, computed on the same
  representation as exact retrieval, with a threshold as the only difference
  \citep{foldiak2003}.
\item \textbf{Generalisation is a subset.} A concept more general than $x$ is
  a \emph{subset} of $x$'s code. Dropping bits is abstracting. This is why
  Borges's Funes, who remembered every instant in full detail and could
  therefore drop nothing, was ``virtually incapable of general, platonic
  ideas'' \citep{borges1942}: no compression, no concepts.
\end{enumerate}

\subsection{The code matrix is a formal context}
\label{sec:code-is-context}

Here is the bridge, stated plainly. Everything in
Section~\ref{sec:fca} is a development of this single sentence:

\begin{keyidea}
A binary stimulus $\times$ neuron matrix and a \emph{formal context}
$\ctx = (G, M, I)$ are the same object. Objects $G$ are the stimuli,
attributes $M$ are the neurons, and the incidence $I$ says ``this neuron fired
for this stimulus''. Everything \fca{} proves about formal contexts is
therefore a theorem about neural codes.
\end{keyidea}

This is not a metaphor invented for this paper; it is a published research
programme. \citet{endres2008} applied \fca{} to recordings from macaque
superior temporal sulcus, building the formal context from spike counts by an
exact Bayesian thresholding procedure and then reading the concept lattice as
a description of what the neural population represents. The lattice exposed a
hierarchical organisation of face representations and evidence for a
product-of-experts style code --- structure that is invisible in the raw
firing rates and that no dimensionality reduction into a Euclidean space would
have expressed, because the structure is \emph{ordinal}, not metric. The
follow-up \citep{endres2010} pushed the same machinery toward semantic
decoding.

\begin{figure}[t]
\centering
\begin{tikzpicture}[font=\small,
  b/.style={draw, rounded corners=3pt, align=center, inner sep=5pt,
            minimum height=1.1cm, text width=3.0cm},
  ar/.style={-{Stealth[length=2.6mm]}, thick},
  lab/.style={font=\scriptsize\itshape, text=odgrey, align=center}]

\node[b, fill=odsand] (brain) at (0,0)
  {\textbf{brain side}\\[2pt]\scriptsize which units\\\scriptsize fire for what};
\node[b, draw=odblue] (matrix) at (4.4,0)
  {\textbf{binary matrix}\\[2pt]\scriptsize stimuli $\times$ units};
\node[b, draw=odblue] (ctx) at (8.8,0)
  {\textbf{formal context}\\[2pt]\scriptsize $(G,M,I)$};
\node[b, fill=odsand] (mind) at (4.4,-2.4)
  {\textbf{mind side}\\[2pt]\scriptsize a hierarchy\\\scriptsize of concepts};
\node[b, draw=odblue] (lat) at (8.8,-2.4)
  {\textbf{concept lattice}\\[2pt]\scriptsize $\Bcl(G,M,I)$};

\draw[ar] (brain) -- node[lab,above]{record\\+ threshold} (matrix);
\draw[ar] (matrix) -- node[lab,above]{rename} (ctx);
\draw[ar] (ctx) -- node[lab,right,xshift=1mm]{close\\(\S\ref{sec:derivation})} (lat);
\draw[ar] (lat) -- node[lab,above]{interpret} (mind);
\draw[ar, dashed, odgrey] (mind) to[bend left=18]
  node[lab,above]{and back} (brain);
\end{tikzpicture}
\caption{The mind--brain round trip that motivates the whole exercise. The
left column is neuroscience, the right column is mathematics, and the two
middle arrows are pure renaming. Because the maps in
Section~\ref{sec:derivation} are invertible, the trip can be made in either
direction: a code determines a concept hierarchy, and a concept hierarchy can
be embedded into a code.}
\label{fig:mind-brain}
\end{figure}

\subsection{Why this is more than a curiosity}

Two directions of travel, both useful.

\paragraph{Code $\to$ concepts (interpretation).} Given a recorded population
code, the concept lattice tells you what the population distinguishes and how
its distinctions nest. This is a tool for neuroscience, and it has the
attractive property of being \emph{non-parametric about geometry}: it assumes
nothing about distances, only about which units fire.

\paragraph{Concepts $\to$ code (design).} Given an ontology you want a system
to hold, you can embed it in a binary code such that subsumption is subset
containment. This is the direction \odag{} cares about, and
Section~\ref{sec:dictionary} shows it is exact: the reflexive ancestor set of
each node \emph{is} such a code, and the graph can be reconstructed from the
codes with nothing lost. Historically the same idea has been rediscovered
several times: as bitvector encodings of type lattices for fast subsumption
checks in programming languages \citep{aitkaci1989}; as faceted
classification in library science, where \citet{ranganathan1933} proposed
describing a work by a \emph{set} of independent facets rather than one
shelf path --- the librarians' independent discovery that the lattice, not
the tree, is the honest shape; as sparse distributed
representations and hyperdimensional computing
\citep{kanerva2009,plate1995}; and, three centuries earlier, as Leibniz's
prime-number notation.

\subsection{A note on Leibniz}
\label{sec:leibniz-note}

Leibniz's scheme deserves a paragraph, because it is the same lattice written
in a different alphabet and it makes the algebra vivid \citep{leibniz1666}.
Assign a distinct prime to each primitive concept, and give a composite
concept the \emph{product} of its constituents' primes. If
$\textsf{animal} = 2$ and $\textsf{rational} = 3$, then $\textsf{human} = 6$.
Now:
\[
  A \text{ subsumes } B \iff \text{the number of } A \text{ divides the number
  of } B .
\]
Divisibility of square-free numbers is exactly set containment on the sets of
prime factors. So Leibniz's ``let us calculate'' and the bitwise-\textsc{and}
of two cone bitmaps are the same operation in different notation, three and a
half centuries apart. What Leibniz lacked was not the algebra but a
\emph{source} for the primitives; \fca{} supplies one (derive them from data),
and \odag{} supplies another (let people assert them, and merge what they
assert).

\subsection{Why not just use vectors?}

A reader who has met machine learning will ask why one would use sets of bits
rather than real-valued embeddings, which are what modern systems actually
use. The honest answer is that they solve different problems, and the
difference is the theme of this paper.

\begin{center}
\small
\begin{tabularx}{\textwidth}{@{}lXX@{}}
\toprule
 & \textbf{Sparse binary code / DAG} & \textbf{Dense real embedding}\\
\midrule
Subsumption & exact set containment, decidable & approximate, thresholded\\
Similarity  & overlap count, exact & cosine distance, learned\\
Query       & intersection: complete and sound & nearest neighbours: neither\\
Novel item  & new bit, no retraining & re-embed, possibly retrain\\
Explanation & the shared bits \emph{are} the reason & post-hoc at best\\
Identity    & a name; stable & a vector; changes when the model does\\
Merging two & union of assertions, converges & incomparable spaces\\
\bottomrule
\end{tabularx}
\end{center}

Embeddings are superb at recall from vague input, which is precisely what an
exact structure cannot do. Discrete structures are superb at exactness,
verifiability and merging, which is precisely what an embedding cannot do.
Section~\ref{sec:agents} argues that a practical agent memory wants both, with
the embedding proposing and the DAG disposing --- and that the interesting
engineering is in the interface between them, not in a winner.

\begin{tryit}
Take the four-stimulus code at the start of this section. Which stimuli
activate a superset of $\{n_2, n_3\}$? (Answer: only \emph{hand}.) Which pairs
of stimuli have overlapping codes? Sketch the containment order on the four
codewords --- you will find that \emph{face A} and \emph{face B} are
incomparable, sharing $\{n_1, n_2\}$, which is the code of a concept that has
no name in the table but plainly ought to. That nameless shared concept is a
formal concept, and finding all of them is what Section~\ref{sec:fca} is
about.
\end{tryit}
